% ============================================================
% Chapter 6: Experimental Design
%   Section: Construction of the Custom Validation Set
%     Subsection: The Tuning and Test Halves
% Included from chapters/06-experimental-design/03-custom-validation-set/custom-validation-set.tex
% Sources: labeller/app.py (split, assign_splits);
% backend/data/evaluation/split.json; docs/decisions/evaluation.md
% ("The split").
% ============================================================

\subsection{The Tuning and Test Halves}
\label{subsec:custom-split}

The system is not trained on the set: it searches the open web when a
claim is checked, and nothing it learns from a labelled claim is stored.
What can still be fitted to the labels is everything the author decides
by looking at results --- a threshold, the trigger of a second search
(Section~\ref{subsec:second-look}), a prompt. Reporting on the same
claims that guided those decisions would inflate the result, so before
any run over the set it was divided once, on 19~October~2026, into a
\textbf{tuning half} (\texttt{train}) and a \textbf{test half}
(\texttt{test}). Every decision is taken on the first; the second is
read once, for the numbers this chapter reports.

\paragraph{Half and half.} The plan first set aside a third for tuning.
Because the pipeline learns nothing from the claims, the risk the split
guards against is limited to the few decisions above, and the author
chose an even split: the tuning half then holds enough disagreements
between system and label to decide on, and the test half still holds
70 claims.

\paragraph{How the halves were drawn.} Three rules, applied by the
labelling tool (\texttt{labeller/app.py split}) rather than by hand:

\begin{itemize}
    \item \textbf{Claims travel with their article.} Claims of one
    article share its evidence, and a synthetic refuted claim
    (Subsection~\ref{subsec:custom-mutations}) is its source claim with
    one element changed and the same evidence. Split apart, tuning on
    one would be tuning on the other; the 139 claims form 59 such
    groups, and a group is never divided.
    \item \textbf{Each half holds its share of every stratum}: of each
    verdict in each language, and of each topic group. Groups are placed
    one at a time, the group holding the rarest stratum first, on the
    side whose targets its claims fall furthest short of --- iterative
    stratification \cite{sechidis2011stratification}, with groups in
    place of single claims. A first draw balanced verdict and language
    only, and put all 21 culture claims in the test half; adding the
    topic group as a second dimension corrected it.
    \item \textbf{Ties are broken by a hash} (SHA-256) of the group's
    first identifier, so no step depends on a choice, and the
    assignment is frozen: a claim keeps its side when it is relabelled,
    and a claim added later joins its article's side, or the side its
    stratum is short of.
\end{itemize}

Table~\ref{tab:custom-split} gives the result. The single
\texttt{MISLEADING} claim fell in the test half, and the two Spanish
\texttt{UNVERIFIED} claims belong to one article and fell in the
tuning half; strata of one or two claims cannot be shared.

\begin{table}[ht]
\centering
\small
\begin{tabular}{@{}lrr@{}}
\toprule
& \textbf{Tuning} & \textbf{Test} \\
\midrule
Claims (articles)        & 69 (28) & 70 (31) \\
\texttt{TRUE}            & 45 & 46 \\
\texttt{PARTIALLY\_TRUE} &  3 &  4 \\
\texttt{MISLEADING}      &  0 &  1 \\
\texttt{UNVERIFIED}      &  6 &  4 \\
\texttt{FALSE} (synthetic) & 15 & 15 \\
\midrule
English / Spanish        & 48 / 21 & 47 / 23 \\
Society / science / environment / culture / health & 11 / 17 / 16 / 10 / 15 & 9 / 16 / 18 / 11 / 16 \\
\bottomrule
\end{tabular}
\caption[The two halves of the custom set]{The two halves of the custom set, frozen on 19 October 2026
(\texttt{backend/data/evaluation/split.json}).}
\label{tab:custom-split}
\end{table}

\paragraph{The re-check sample.} The fifth of the test half that is
labelled a second time (Subsection~\ref{subsec:custom-review}) is
drawn from it by the same hash, and recorded with the split: 14 claims.
The same 14 claims are also run through the system a second time
(Subsection~\ref{subsec:consistency}), so the consistency of the labels
and the consistency of the system are measured on one sample.
